\chapter{Arithmetik und die natürliche Ordnung}
\label{b10guide:ch:arithmetic}

Alle Zahlenbuchstaben in diesem Kapitel bezeichnen natürliche
Zahlen. Es gilt $1=S(0)$. Die Rechenoperationen werden durch die
Rekursion auf $\mathbb N$ festgelegt; ihre Gesetze sind Folgerungen
der Peano-Eigenschaften.

\section{Addition und Multiplikation}

\begin{definition}[Rekursive Festlegung der Addition]
\label{b10guide:def:addition}
Die Addition ist die eindeutig bestimmte Operation auf $\mathbb N$ mit
\[
 a+0=a,\qquad a+S(b)=S(a+b).
\]
Insbesondere gehören alle Summen wieder zu $\mathbb N$.
\end{definition}

\begin{theorem}[Grundgesetze der Addition]
\label{b10guide:thm:addition-laws}
Für natürliche Zahlen gelten
\[
 \begin{aligned}
 a+0&=0+a=a, & a+1&=S(a),\\
 a+b&=b+a, & (a+b)+c&=a+(b+c),\\
 S(a)+b&=S(a+b), & a+S(b)&=S(a+b).
 \end{aligned}
\]
Außerdem gelten die Kürzungs- und Nullsummenregeln
\[
 a+c=b+c\ \Longleftrightarrow\ a=b,
 \qquad
 a+b=0\ \Longleftrightarrow\ a=b=0.
\]
\end{theorem}

\begin{definition}[Rekursive Festlegung der Multiplikation]
\label{b10guide:def:multiplication}
Die Multiplikation ist die eindeutig bestimmte Operation auf
$\mathbb N$ mit
\[
 a\cdot0=0,\qquad a\cdot S(b)=a\cdot b+a.
\]
Alle Produkte gehören zu $\mathbb N$. Der Punkt kann in Produkten
weggelassen werden.
\end{definition}

\begin{theorem}[Grundgesetze der Multiplikation]
\label{b10guide:thm:multiplication-laws}
Für natürliche Zahlen gelten
\[
 \begin{aligned}
 a\cdot0&=0\cdot a=0, & a\cdot1&=1\cdot a=a,\\
 ab&=ba, & (ab)c&=a(bc),\\
 a(b+c)&=ab+ac, & (a+b)c&=ac+bc.
 \end{aligned}
\]
\end{theorem}

\begin{theorem}[Kürzung, Nullprodukte und Produktgleichheit mit eins]
\label{b10guide:thm:product-cancellation}
Es gelten
\[
 ab=0\ \Longleftrightarrow\ (a=0\ \lor\ b=0),\qquad
 ab=1\ \Longleftrightarrow\ a=b=1.
\]
Für $c\ne0$ gilt
\[
 ac=bc\ \Longleftrightarrow\ a=b.
\]
Die Voraussetzung $c\ne0$ gehört zur Kürzungsregel.
\end{theorem}

\section{Vergleich und natürlicher Abstand}

\begin{definition}[Natürliche Ordnung]
\label{b10guide:def:natural-order}
Für $m,n\in\mathbb N$ bedeuten
\[
 m\le n\quad\Longleftrightarrow\quad
 \exists k\in\mathbb N\ (m+k=n),
 \qquad
 m<n\quad\Longleftrightarrow\quad(m\le n\ \land\ m\ne n).
\]
Die Schreibweisen $n\ge m$ und $n>m$ kehren die Argumente um.
Sie beziehen sich hier ausschließlich auf die natürlichen Zahlen.
\end{definition}

\begin{theorem}[Eindeutiger Abstand und Vergleich]
\label{b10guide:thm:natural-distance}
Ist $m\le n$, so gibt es genau ein $k\in\mathbb N$ mit $m+k=n$.
Ferner gilt
\[
 m<n\ \Longleftrightarrow\
 \exists!k\in\mathbb N\ (n=m+k+1).
\]
Für je zwei natürliche Zahlen gilt genau einer der drei Fälle
$m<n$, $m=n$ oder $n<m$.
\end{theorem}

\begin{definition}[Differenz bei nichtnegativem Abstand]
\label{b10guide:def:natural-difference}
Für $m\le n$ bezeichnet $n-m$ die eindeutig bestimmte natürliche
Zahl $k$ mit $m+k=n$. Die Schreibweise $n-m$ wird in diesem
Leitfaden nur unter der Voraussetzung $m\le n$ verwendet.
\end{definition}

\begin{theorem}[Rechnen mit zulässigen Differenzen]
\label{b10guide:thm:natural-difference}
Für $m\le n$ gelten
\[
 m+(n-m)=n,\qquad (n+c)-(m+c)=n-m.
\]
Für $\ell\le m\le n$ gelten
\[
 (n-m)+(m-\ell)=n-\ell,\qquad
 (n-\ell)-(m-\ell)=n-m.
\]
Alle in diesen Formeln auftretenden Differenzen sind im Sinn
von Definition~\ref{b10guide:def:natural-difference} zulässig.
\end{theorem}

\begin{theorem}[Null, Nachfolger und Diskretheit]
\label{b10guide:thm:discrete-order}
Für $m,n\in\mathbb N$ gelten
\[
 \begin{gathered}
 0\le n,\qquad n\ne0\ \Longleftrightarrow\ 1\le n,\\
 m<n\ \Longleftrightarrow\ m+1\le n,
 \qquad m<n+1\ \Longleftrightarrow\ m\le n.
 \end{gathered}
\]
Zwischen $n$ und $n+1$ liegt keine natürliche Zahl. Jede positive
natürliche Zahl besitzt den unmittelbaren Vorgänger $n-1$, und
zu jeder natürlichen Zahl $n$ gibt es die größere Zahl $n+1$.
\end{theorem}

\begin{theorem}[Vergleich von Summen und Produkten]
\label{b10guide:thm:arithmetic-order}
Für natürliche Zahlen gelten
\[
 \begin{aligned}
 a\le b&\ \Longleftrightarrow\ a+c\le b+c,\\
 a<b&\ \Longleftrightarrow\ a+c<b+c.
 \end{aligned}
\]
Aus $a\le b$ und $c\le d$ folgen $a+c\le b+d$ und $ac\le bd$.
Für $c>0$ gelten außerdem
\[
 a\le b\ \Longleftrightarrow\ ac\le bc,
 \qquad
 a<b\ \Longleftrightarrow\ ac<bc.
\]
\end{theorem}

\begin{theorem}[Größenvergleich mit Summanden und Faktoren]
\label{b10guide:thm:arithmetic-bounds}
Es gelten $a\le a+b$ und
\[
 a+b=a\ \Longleftrightarrow\ b=0.
\]
Für $b\ge1$ gilt $a\le ab$. Für $a>0$ gilt
\[
 ab=a\ \Longleftrightarrow\ b=1.
\]
Insbesondere ist für $a>0$ und $b>1$ stets $a<ab$.
\end{theorem}

\section{Natürliche Anfangsabschnitte und ihre Mengendarstellung}

Die folgenden Regeln betreffen die natürlichen Anfangsabschnitte
$I_n=\{k\in\mathbb N\mid k<n\}$ aus
Definition~\ref{b10guide:def:recursion-history}.
Sie entsprechen den Abschnitten über Peano-Anfangsabschnitte und
Blockzerlegung im Hauptband \cite{project}.

\begin{theorem}[Nachfolger und Vergleich natürlicher Anfangsabschnitte]
\label{b10guide:thm:initial-segments}
Für $m,n\in\mathbb N$ gilt
\[
 \begin{gathered}
 I_0=\varnothing,\qquad I_{n+1}=I_n\cup\{n\},\qquad n\notin I_n,\\
 I_m\subseteq I_n\ \Longleftrightarrow\ m\le n,\\
 I_m\subsetneq I_n\ \Longleftrightarrow\ m<n,\qquad
 I_m=I_n\ \Longleftrightarrow\ m=n.
 \end{gathered}
\]
\end{theorem}

\begin{theorem}[Disjunkte Blockzerlegung eines natürlichen Anfangsabschnitts]
\label{b10guide:thm:initial-blocks}
Für $m,n\in\mathbb N$ gilt
\[
 \begin{gathered}
 I_{m+n}=I_m\cup\{m+j\mid j\in I_n\},\\
 I_m\cap\{m+j\mid j\in I_n\}=\varnothing.
 \end{gathered}
\]
Jede Zahl $i\in I_{m+n}$ liegt entweder in $I_m$ oder besitzt
genau eine Darstellung $i=m+j$ mit $j\in I_n$.
Dies schließt die leeren Blöcke bei $m=0$ oder $n=0$ ein.
\end{theorem}

\begin{theorem}[Mitgliedschaft und Inklusion in der von-Neumann-Realisierung]
\label{b10guide:thm:von-neumann-order}
Im Peano-System $(\omega,\varnothing,S)$ mit $S(n)=n\cup\{n\}$
aus Satz~\ref{b10guide:thm:von-neumann} gilt für $m,n\in\omega$
\[
 \begin{gathered}
 m\in n\ \Longleftrightarrow\ m<n\ \Longleftrightarrow\ m\subsetneq n,\\
 m\subseteq n\ \Longleftrightarrow\ m\le n,\qquad
 n=\{k\in\omega\mid k<n\}.
 \end{gathered}
\]
Dabei sind $<$ und $\le$ die durch die Addition dieses Peano-Systems
definierten Vergleiche. Wird $\mathbb N=\omega$ gewählt, gilt somit
$I_n=n$. Diese Mengengleichheit wird für ein beliebig gewähltes
Peano-System nicht vorausgesetzt.
\end{theorem}

\begin{theorem}[Größte Zahl in einer beschränkten natürlichen Teilmenge]
\label{b10guide:thm:bounded-maximum}
Sei $\varnothing\ne U\subseteq\mathbb N$. Gibt es ein $b\in\mathbb N$
mit $u\le b$ für alle $u\in U$, so gibt es genau ein $M\in U$ mit
\[
 \forall u\in U\;(u\le M).
\]
\end{theorem}
